\documentclass[11pt, a4paper]{article}

\usepackage[T1]{fontenc}
\usepackage[utf8]{inputenc}
\usepackage{tgpagella}
\usepackage{microtype}
\usepackage[spanish, es-noshorthands]{babel}
\addto\captionsspanish{\renewcommand{\tablename}{Tabla}}

\usepackage{amsmath, amssymb, amsfonts}
\usepackage{graphicx}
\usepackage{booktabs}
\usepackage[most]{tcolorbox}
\usepackage[margin=2.5cm]{geometry}
\usepackage{fancyhdr}

\pagestyle{fancy}
\fancyhf{}
\setlength{\headheight}{16pt}
\lhead{\textbf{Ing. Mecatrónica}}
\chead{Examen EC1 - Versión B}
\rhead{Robótica (IMT-342)}
\lfoot{Prof: M.Sc. Ing. Bernardo Quiroga Turdera}
\rfoot{Página \thepage}
\renewcommand{\headrulewidth}{0.4pt}
\renewcommand{\footrulewidth}{0.4pt}

\definecolor{ucbblue}{RGB}{0, 51, 102}

\begin{document}

\begin{center}
    \Large\textbf{\textcolor{ucbblue}{EC1 — Evaluación Diferida (Versión B)}}\\[0.2cm]
\end{center}

\textbf{Nombre:} \rule{8cm}{0.15mm} \hfill \textbf{Fecha:} \rule{4cm}{0.15mm}

\vspace{0.4cm}
\begin{tcolorbox}[colback=gray!5, colframe=gray!50, title=\textbf{Condiciones de la Evaluación[cite: 8, 9]}]
    \begin{itemize}
        \item Examen escrito individual. No se permite el uso de calculadora ni formulario.
        \item Muestre su razonamiento y transformaciones intermedias; respuestas finales sin respaldo tendrán crédito limitado.
        \item Utilice la convención D-H \textbf{estándar} estudiada en clase.
        \item A menos que se especifique, las distancias tienen unidades arbitrarias consistentes y los ángulos están en radianes.
    \end{itemize}
\end{tcolorbox}

\section*{Problema 1 — Recuperación de Transformación (20 pts)[cite: 8]}
Se conocen las siguientes matrices de transformación:
\begin{equation*}
    {}^{0}T_2 = \begin{bmatrix} 0 & -1 & 0 & 1 \\ 1 & 0 & 0 & 2 \\ 0 & 0 & 1 & 0 \\ 0 & 0 & 0 & 1 \end{bmatrix}, \qquad
    {}^{1}T_2 = \begin{bmatrix} 1 & 0 & 0 & 1 \\ 0 & 1 & 0 & 0 \\ 0 & 0 & 1 & 0 \\ 0 & 0 & 0 & 1 \end{bmatrix}
\end{equation*}

\textbf{a) [6 pts]} Escriba la relación de encadenamiento de marcos que debe cumplirse entre ${}^{0}T_1$, ${}^{1}T_2$ y ${}^{0}T_2$.

\textbf{b) [8 pts]} Determine analíticamente ${}^{0}T_1$.

\textbf{c) [6 pts]} Explique físicamente por qué la traslación contenida en ${}^{1}T_2$ no puede simplemente restarse componente a componente de la traslación de ${}^{0}T_2$ sin considerar la orientación.

\section*{Problema 2 — Detección y Corrección D-H (25 pts)[cite: 8]}
Una articulación revoluta se describe mediante la siguiente fila D-H \textbf{estándar}:
\begin{equation*}
    \theta_i = q_i, \qquad d_i = 1, \qquad a_i = 2, \qquad \alpha_i = 0
\end{equation*}
Un estudiante propone la siguiente matriz de transformación:
\begin{equation*}
    \tilde{T} = \begin{bmatrix} \cos q_i & -\sin q_i & 0 & 2 \\ \sin q_i & \cos q_i & 0 & 1 \\ 0 & 0 & 1 & 0 \\ 0 & 0 & 0 & 1 \end{bmatrix}
\end{equation*}

\textbf{a) [7 pts]} Identifique al menos \textbf{dos} errores conceptuales en la matriz propuesta por el estudiante.

\textbf{b) [10 pts]} Construya la matriz de transformación homogénea correcta ${}^{i-1}T_i(q_i)$.

\textbf{c) [4 pts]} Evalúe la matriz correcta en $q_i = \pi/2$.

\textbf{d) [4 pts]} Indique cuál parámetro D-H es la variable articular y por qué.

\section*{Problema 3 — FK de Representación Mixta (45 pts)[cite: 8]}
Un manipulador serial 3R utiliza convención D-H estándar. En lugar de darle la tabla completa, ya se conocen las dos primeras transformaciones:
\begin{equation*}
    {}^{0}T_1 = \begin{bmatrix} 0 & -1 & 0 & 0 \\ 1 & 0 & 0 & 1 \\ 0 & 0 & 1 & 0 \\ 0 & 0 & 0 & 1 \end{bmatrix}, \qquad
    {}^{1}T_2 = \begin{bmatrix} 1 & 0 & 0 & 1 \\ 0 & 0 & -1 & 0 \\ 0 & 1 & 0 & 0 \\ 0 & 0 & 0 & 1 \end{bmatrix}
\end{equation*}
La tercera articulación es revoluta y tiene los parámetros:
\begin{equation*}
    \theta_3 = q_3, \qquad d_3 = 0, \qquad a_3 = 1, \qquad \alpha_3 = 0 \qquad \text{con} \quad q_3 = \pi
\end{equation*}

\textbf{a) [8 pts]} Construya ${}^{2}T_3$.

\textbf{b) [5 pts]} Escriba la cadena de transformación correcta para ${}^{0}T_3$.

\textbf{c) [16 pts]} Calcule numéricamente ${}^{0}T_3$.

\textbf{d) [8 pts]} Extraiga ${}^{0}p_3$ y ${}^{0}R_3$, y explique qué representa cada uno.

\textbf{e) [8 pts]} Dé una verificación estructural y una verificación física/cinemática de su resultado.

\section*{Problema 4 — Verificación de Estructura (10 pts)[cite: 8]}
Un estudiante reporta que el resultado de cinemática directa de un manipulador planar 2R es:
\begin{equation*}
    \tilde{T} = \begin{bmatrix} 0 & -1 & 0 & 1 \\ 1 & 0 & 0 & 1 \\ 0 & 0 & 1 & 0 \\ 0 & 0 & 1 & 1 \end{bmatrix}
\end{equation*}

\textbf{a) [5 pts]} Sin recalcular la cinemática directa, identifique una razón \textbf{estructural} por la cual $\tilde{T}$ no puede ser una transformación homogénea de cuerpo rígido.

\textbf{b) [5 pts]} El mismo estudiante argumenta: \textit{"La posición $[1, 1, 0]^T$ parece físicamente plausible, por lo tanto la transformación completa debe ser correcta."} Explique por qué ese argumento es insuficiente y mencione \textbf{una verificación independiente adicional} que debería realizarse.

\end{document}
